\documentclass[12pt,a4paper]{article}
\usepackage{fontspec}
\newfontfamily\cjkfont{Droid Sans Fallback}
\newcommand{\zh}[1]{{\cjkfont #1}}
\usepackage[spanish]{babel}
\usepackage{geometry}
\geometry{margin=1in}
\usepackage{hyperref}
\usepackage{setspace}
\onehalfspacing
\usepackage{parskip}

\begin{document}

% --- Cover page ---
\begin{titlepage}
\centering
\vspace*{3cm}

{\LARGE\bfseries Tecnología, política y el docente que necesitamos\par}

\vspace{2cm}

{\large Hesus García Cobos\par}

\vspace{1cm}

{\large Maestría en Pedagogía\par}
{\large Universidad Popular Autónoma del Estado de Puebla\par}

\vspace{1cm}

{\large Profesora: Socorro Guevara Salazar\par}

\vspace{1cm}

{\large 19 de marzo de 2026\par}

\vfill
\end{titlepage}

% --- Reflection body ---

Justo estoy reflexionando sobre esto porque tuve la oportunidad de viajar a China en 2024 y porque antes estudié en Canadá. Vi dos sistemas educativos muy diferentes del mexicano, y los dos me dejaron pensando en lo mismo: la tecnología en el aula no es cuestión de tener la herramienta, es cuestión de cómo la usas y qué hay detrás.

En China lo que me sorprendió fue la disciplina. No solo de los alumnos, del sistema entero. El gobierno publicó un plan nacional de IA en 2017 (Consejo de Estado, 2017), y en menos de ocho años logró que todas las primarias y secundarias del país ofrezcan mínimo 8 horas anuales de instrucción en inteligencia artificial. Ronghuai Huang (\zh{黄荣怀}) de la Universidad Normal de Pekín lleva años trabajando en lo que llama ``educación inteligente'': la idea de que la tecnología debe rediseñar cómo se aprende, no solo sustituir al pizarrón (Huang, 2025). Escogí a Huang porque es de los pocos investigadores que trabaja directo con el gobierno chino y a la vez publica para audiencias internacionales, entonces puedes ver la política y la teoría juntas. China también posicionó a Qian Tang (\zh{唐虔}) como subdirector general de educación en la UNESCO, conectando su modelo educativo con la agenda global. Todo se mueve con dirección clara desde arriba.

En Canadá la historia es diferente. Lo que vi ahí fue infraestructura. Acceso a internet que funciona, laboratorios equipados, profesores que experimentan con herramientas digitales porque las tienen disponibles. No hay un mandato central como en China, el sistema te da las condiciones y tú decides qué hacer con ellas.

Y en México, pues ni lo uno ni lo otro. Ni la disciplina centralizada de China ni la infraestructura de Canadá. Le pedimos competencia digital al docente sin darle internet estable, y le exigimos innovación sin tiempo ni espacio para experimentar.

El modelo del traje docente describe cuatro piezas: empatía, competencia digital, reflexión pedagógica y diseño instruccional. Esas cuatro piezas me sirven para ordenar lo que vi en estos dos países y lo que veo en México.

\section*{¿Qué áreas resultan más desafiantes para el docente actual?}

La más difícil es usar tecnología con sentido pedagógico. No usar una aplicación porque sí, sino que la herramienta resuelva un problema de aprendizaje que sin ella no se resuelve igual. Esto es lo que Kalantzis y Cope (2012) describen como el cambio del docente de transmisor a diseñador de experiencias de aprendizaje. Me gusta esta postura porque no dice ``usa más tecnología'', dice ``piensa diferente sobre lo que haces en el aula''. El traje docente lo plantea con los marcos de conocimiento tecnológico pedagógico del contenido y el modelo de sustitución, aumento, modificación y redefinición: la tecnología no sustituye al maestro, lo obliga a diseñar distinto.

La otra área difícil es mantener autonomía cuando el sistema te dice exactamente qué enseñar y cómo. Me respaldo en Apple (2018), que lleva décadas argumentando que la centralización curricular convierte al profesor en ejecutor de lo que otros decidieron. Escogí a Apple porque su crítica aplica perfecto al caso chino: China implementó IA en todas sus escuelas, pero ¿cuánta libertad tiene el maestro de una comunidad rural en Guizhou para adaptar esas directivas a su contexto? La propia UNESCO advierte que las políticas de IA deben respetar la autonomía del docente (Miao et al., 2021).

\section*{¿Qué componentes del traje requieren mayor desarrollo en el contexto educativo?}

El digital, pero no aislado. La pieza digital sin reflexión pedagógica produce cursos llenos de aplicaciones y vacíos de aprendizaje. Lo que falta desarrollar es la articulación entre lo digital y lo reflexivo: la capacidad del docente de experimentar con tecnología, evaluar si funcionó y ajustar. Ni China ni Canadá tienen esto resuelto del todo. China tiene la escala pero le falta autonomía. Canadá tiene la infraestructura pero no siempre la dirección. México necesita las dos cosas y la conexión entre ellas.

\section*{¿Cómo pueden las instituciones educativas apoyar la construcción de este nuevo perfil docente?}

Con tres cosas concretas. Primera, formación continua que sirva, no cursos donde te sientan a ver diapositivas sobre herramientas que nunca vas a usar. Segunda, acceso a infraestructura: no se puede pedir competencia digital si el profesor no tiene ni internet estable. Tercera, espacios de experimentación, como propone la UNESCO (2022): tiempo institucional para probar cosas sin que cada intento cuente como evaluación de desempeño. Escogí este documento de la UNESCO porque es de los pocos que habla de lo que las instituciones deben cambiar, no solo de lo que el docente debe aprender.

\section*{Cierre}

El traje docente no viene hecho. No puedes copiarlo de China ni de Canadá. Las cuatro piezas se construyen en la práctica, con apoyo institucional pero con criterio propio. China me enseña que sí se puede mover un sistema educativo entero en pocos años. Canadá me enseña que la infraestructura importa más de lo que creemos. Y México me recuerda que sin condiciones básicas, cualquier modelo se queda en discurso.

% --- References ---
\newpage
\begin{center}
{\large\bfseries Referencias}
\end{center}
\vspace{0.5cm}

\noindent Apple, M. W. (2018). \textit{Ideology and curriculum} (4a ed.). Routledge.

\vspace{0.4cm}
\noindent Consejo de Estado de la República Popular China. (2017). \textit{Plan de desarrollo de inteligencia artificial de nueva generación} (\zh{国务院关于印发新一代人工智能发展规划的通知}).

\vspace{0.4cm}
\noindent Huang, R. (2025). \textit{Smart education: Pathways toward education 2050 (Abridged Edition)}. Smart Learning Institute, Universidad Normal de Pekín.

\vspace{0.4cm}
\noindent Kalantzis, M., y Cope, B. (2012). \textit{New learning: Elements of a science of education} (2a ed.). Cambridge University Press.

\vspace{0.4cm}
\noindent Miao, F., Holmes, W., Huang, R., y Zhang, H. (2021). \textit{AI and education: Guidance for policy-makers}. UNESCO.

\vspace{0.4cm}
\noindent UNESCO. (2022). \textit{Reimaginar juntos nuestros futuros: Un nuevo contrato social para la educación}. UNESCO.

\end{document}
